\section{Introduction}\label{sec:intro}

Electro-optical cameras are the cheapest sensor on a vessel that can see small, non-cooperative objects: buoys, kayaks, drifting debris, birds on the water and small craft that carry no AIS transponder. Detecting these objects reliably at range is a precondition for collision avoidance on autonomous surface vessels and for decision support on crewed ships. Range makes the task hard: at the distances that matter for manoeuvring, most objects occupy a few pixels of a high-resolution frame, and the conventional practice of resizing a frame to a fixed detector input removes much of that information.

SeaClips \citep{denk2026seaclips} is a recent video benchmark for this problem. It provides 74 annotated videos (47 for training, 9 for validation and 18 for testing) with seven classes, and reports nine baselines evaluated with COCO-style average precision at an intersection-over-union of 0.3 (\mapthree{}), averaged over three seeds. The strongest published results are TransVOD++ with a Swin-B backbone ($31.95 \pm 3.28$) and Deformable DETR with Swin-B ($31.28 \pm 3.72$); Deformable DETR with ResNet-50 reaches $22.86 \pm 0.91$.

Two properties of SeaClips make it hard to know, during development, whether a design choice will help on the test split. First, most objects are very small: at the benchmark's $640 \times 640$ input, the median training box of the ANIMAL class is under 5 pixels across. Second, the validation split represents some test classes poorly or not at all: it contains no LEISURE\_VEHICLE boxes, and all of its ANIMAL boxes come from one video. Our central question is therefore not only whether a detector can score higher on TEST, but which development decisions, made on TRAIN and VAL, transfer to unseen test videos.

We answer it with a pre-registered programme of fifteen gates. For each gate, the hypothesis, the pass rule and the consequence of failure were written down and pinned by file hash before the run. The final system and two pre-declared variants were evaluated once on TEST. The label file had been hash-locked without being parsed, and the evaluation tool refused to run a second time.

The contributions are as follows.
\begin{enumerate}
\item An evaluation of resolution-preserving tile training and 17-tile inference with a ConvNeXt-B Deformable-DETR on SeaClips over three seeds. It is compared with the published baselines and with our own same-pipeline reproduction of the R50 baseline, and reports latency next to accuracy (Section~\ref{sec:results}).
\item A complete account of fourteen negative or partial results. These include a validation-based per-class routing decision that transferred negatively to TEST (Section~\ref{sec:negative}).
\item Five measured properties of SeaClips that bound what its validation split can show about its test split, with recommendations for future evaluations (Section~\ref{sec:dataset}).
\item Released code and pre-registration records. In line with the dataset licence, no images, annotations, predictions or weights are redistributed (Section~\ref{sec:statements}).
\end{enumerate}

We do not claim that the approach outperforms TransVOD++ or Deformable-DETR Swin-B: under the benchmark's own test, the differences in mean are not significant.
